\section{Related Work}\label{sec:related}
The Introduction linked strategic incentives, message delivery, and uncertainty
within a self-hosted simulated world. To situate that design, this section
first examines how agent-based models are specified, and then reviews
game-theoretic and financial environments, social simulation, and probabilistic
decision evaluation. Throughout, comparisons concern modeling scope and
preserved evidence rather than a performance ranking of systems.

\subsection{Agent-Based Simulation and Model Specification}
To make interacting actors a reproducible object of study, agent-based research
requires an explicit description of the environment and its rules. In
particular, the Overview, Design concepts, and Details (\textit{ODD}) protocol
standardizes descriptions of entities, initialization, schedules, inputs, and
submodels \cite{grimm2020}. Complementing specification standards,
Gal\'an et al.\ distinguish conceptual modeling assumptions from artifacts
introduced by software implementation \cite{galan2009}. Together, these works
motivate separating verification---whether code executes stated rules
accurately---from validation of whether those rules represent real markets.

Building on that distinction, \textit{Boyun} specifies actors,
resources, observation boundaries, actions, and settlement rules before a run
begins. Preserving these inputs opens assumptions to review and enables exact
replay checks. However, an explicit specification alone cannot show that
assumed incentives or actor choices are realistic, which requires examining
both the economic environment and the behavioral model.

\subsection{Game-Theoretic and Economic Environments}
On the strategic side, \textit{game theory} links actors' choices to
outcomes that depend on what others do. \textit{Nash equilibrium} provides a
classic benchmark in which no actor gains from a unilateral deviation
\cite{nash1950}, and \textit{OpenSpiel} operationalizes finite games through
explicit states, legal actions, observations, transitions, and returns
\cite{lanctot2019}. \textit{Boyun} adopts explicit game specification to
construct decision worlds; by contrast, it neither incorporates
\textit{OpenSpiel}'s equilibrium solvers nor assumes that its actor decision
model plays an equilibrium.

Beyond abstract games, economic simulation environments enforce
incentives through allocation and accounting rules. For instance, the
\textit{AI Economist} couples economic agents with a tax-policy planner,
separating individual utility from a designer's objective \cite{zheng2020},
while \textit{EconAgent} places language-model agents inside numerical
macroeconomic rules for labor and consumption \cite{li2024econagent}.
Similarly, \textit{Boyun} separates actor preference weights,
action selection, and engine-computed consequences, though it targets bounded
financial and protocol decisions rather than macroeconomic policy learning.

When questions turn specifically to trading and decentralized credit,
settlement rules must reflect venue mechanics. \textit{ABIDES} (Agent-Based
Interactive Discrete Event Simulation) models trading agents communicating with
an order-book exchange over configurable network latency \cite{byrd2019}.
On the other hand, decentralized exchanges commonly replace order books
with \textit{automated market makers} (AMMs), whose reserve formulas govern
pricing, fees, and slippage \cite{angeris2021,xu2023amm}, while decentralized
lending pools price credit through utilization-based interest curves and
collateral liquidations \cite{gudgeon2020}. These precedents show why a
narrative choice or participation score cannot substitute for a settled swap,
loan, or cash balance; \textit{Boyun} therefore pairs its ledgers with an
explicit social-information layer.

\subsection{Social Graphs and \rev{Agent-Swarm} Simulation}
Turning to that social layer, language-model simulations study how
\rev{an \textit{agent swarm} observes and reacts to connected peers across a
\textit{social graph}}. \textit{Generative Agents} combines memory,
reflection, and planning to produce contextual behavior \cite{park2023},
whereas \textit{OASIS} provides large-scale social-platform environments with
posting actions and recommendation feeds \cite{yang2025}. Both frameworks
motivate treating observation access and interaction history as explicit parts
of the world rather than assuming shared global knowledge \rev{across the swarm}.

\rev{Directly inspiring \textit{Boyun},} \textit{MiroFish} \rev{pioneered an}
end-to-end workflow from world construction to \rev{agent-swarm} simulation and
structured reporting \cite{mirofish}. However, whereas \textit{MiroFish} runs
open-ended social discourse on \textit{OASIS}, \textit{Boyun} \rev{models actors
as a \textit{swarm of agents} communicating over an explicit \textit{social
graph} (implemented as a directed information graph) while executing} bounded
actions over deterministic strategic, market, and lending ledgers.
Section~\ref{sec:res-mirofish} compares the two systems on the same launch
brief, documenting their complementary outputs rather than claiming a
performance ranking.

Even when social and financial rules are explicit, plausible prose or
contextual responsiveness does not guarantee incentive-consistent choices.
Controlled repeated-game studies with language models test how models respond
to explicit payoffs and opponent histories \cite{akata2025}. Accordingly,
evaluating \textit{Boyun} requires separating rule-enforced consequences from
model-selected behavior and distinguishing option scores from outcome
probabilities.

\subsection{Probabilistic Decisions and Statistical Evaluation}
To obtain structured choices without free-text parsing, typed decision
models score a finite menu of options given a state \cite{typesafeintro,typesafeapi}.
Because financial scenarios often contain confidential strategy, open-weight
implementations are particularly relevant: \textit{Laya} fine-tunes a
421-million-parameter encoder to return normalized option scores on a local CPU
\cite{layamodel}, and \textit{Kev} trains compatible open-weight variants on
Qwen backbones \cite{kevrepo}, whereas hosted services such as \textit{Jev}
process requests remotely \cite{typesafeapi}. \textit{Boyun} uses local
\textit{Laya} inference in its experiments while keeping feasibility checks and
accounting in code. At the same time, \textit{Laya}'s model card notes
that its training workflows exclude financial markets and that its confidence
values are uncalibrated \cite{layamodel}, so returned option scores must never
be read as empirical market odds.

By contrast, reasoning about external event uncertainty requires explicit
hypotheses, priors, and likelihoods via \textit{Bayesian updating}
\cite{pishronik2014}. Repeated sampling can explore consequences under declared
priors, yet it cannot supply missing likelihoods or correct a biased behavioral
model. Ultimately, any probabilistic forecast claim must be evaluated against
resolved outcomes using strictly proper scoring rules such as the
\textit{Brier} and \textit{logarithmic scores} \cite{brier1950,gneiting2007}.

Taken together, prior work motivates combining explicit game incentives,
directed message delivery, self-hosted typed actor decisions, and ledger-based
settlement, while maintaining a strict boundary between verified mechanics and
unvalidated market forecasts. The next section formalizes how a world is
constructed within that boundary.
